\documentclass[10pt]{beamer}
\usepackage{tikz}
\usetikzlibrary{arrows.meta}
\usetikzlibrary{fit}

\usetikzlibrary{arrows,positioning, calc,lindenmayersystems,decorations.pathmorphing,intersections}
\tikzstyle{resource}= [draw,minimum size=16pt,inner sep=0pt]
\tikzstyle{process} = [draw,minimum size=16pt,inner sep=0pt,circle]
\tikzstyle{allocated} = [->,thick,arrows={-latex}]
\tikzstyle{requested} = [<-,thick,arrows={latex-}, dashed]

\usepackage{tikz}



\usetikzlibrary{positioning}
\newcommand{\MonetaryLevel}{Monetary level}
\newcommand{\RealLevel}{Real level}
\newcommand{\Firms}{Firms}
\newcommand{\Households}{Households}
\newcommand{\Banks}{Banks}
\newcommand{\Commodities}{Commodities}
\newcommand{\LaborPower}{Labor power}
\newcommand{\Wages}{Wages}
\newcommand{\Consumption}{Consumption}
\newcommand{\Credits}{Credits}
\newcommand{\Withdrawals}{Withdrawals}
\newcommand{\Deposits}{Deposits}
\newcommand{\Repayments}{Repayments}

\newcommand{\yslant}{0.5}
\newcommand{\xslant}{-0.6}



\newcommand{\VALID}		{\mathsf{valid}}
\newcommand{\INVALID}	{\mathsf{invalid}}
\newcommand{\UNKNOWN}	{\mathsf{unknown}}
\newcommand{\FAILED}	{\mathsf{failed}}
\newcommand{\prover}	{\mathrm{prover}}
\newcommand{\INTER}[1] 	{\llbracket #1 \rrbracket}

\providecommand{\keywords}[1]{\textbf{{Keyword:}} #1}

\usepackage{bsymb}
\newcommand{\STATE}		{\operatorname{\Omega}}
\newcommand{\fst}		{\operatorname{fst}}
\newcommand{\lst}		{\operatorname{lst}}
\newcommand{\elm}		{\operatorname{elm}}

\newcommand{\FN}[1] 	{{\color{dblue}{\mathrm{#1}}}}
\newcommand{\SET}[1] 	{\mathrm{\textsc{#1}}}
\newcommand{\OP}[1] 	{\mathbf{#1}}
\newcommand{\ALSS}		{\prec}
\newcommand{\ALSE}		{\preccurlyeq}

\newcommand{\CP}[1] 	{{\color{black}{\mathsf{#1}}}}
\newcommand{\LG}[1] 	{{\color{black}{#1}}}

\newcommand{\CLEAR} 	{{\color{dgreen}{\mathsf{clear}}}}
\newcommand{\NORMAL} 	{{\color{dgreen}{\mathsf{normal}}}}
\newcommand{\REVERSE} 	{{\color{dgreen}{\mathsf{reverse}}}}
\newcommand{\OCC} 		{{\color{dgreen}{\mathsf{occupied}}}}
\newcommand{\LOCKED} 	{{\color{dgreen}{\mathsf{locked}}}}

\newcommand{\PART}[1] 	{{\color{dgreen}{\mathrm{#1}}}}
\newcommand{\ELM}[1] 	{{\color{dblue}{\mathrm{#1}}}}



\newcommand{\NODE}		{\CP{N}}
\newcommand{\TRACK}		{\CP{T}}
\newcommand{\ROUTE}		{\CP{R}}
\newcommand{\LINE}		{\CP{L}}
\newcommand{\LINEA}		{\CP{L_A}}
\newcommand{\POINT}		{\CP{P}}
\newcommand{\AMBIT}		{\CP{A}}
\newcommand{\SECTION}	{\CP{S}}
\newcommand{\JUNCTION}	{\CP{J}}
\newcommand{\RRES}		{\FN{rr}}
\newcommand{\PRES}		{\FN{pr}}
\newcommand{\RREQ}		{\FN{rr}}
\newcommand{\XRULE}		{\FN{rule}}
\newcommand{\RINT}[1]	{\left \llbracket \begin{array}{@{}l@{}} #1 \end{array} \right \rrbracket}

\newcommand{\TLINE}		{\FN{\ell}}
\newcommand{\FRONT}		{\FN{hd}}	
\newcommand{\TAIL}		{\FN{tl}}
\newcommand{\TRES}		{\FN{lk}}
\newcommand{\RSET}		{\FN{rs}}
\newcommand{\STATEP}	{\FN{p}}
\newcommand{\STATER}	{\FN{r}}

\newcommand{\FROUTE}	{\FN{rhd}}
\newcommand{\TROUTE}	{\FN{rtl}}
\newcommand{\TRAN}		{\rightarrow}	
\newcommand{\TRAIN}		{\SET{tr}}	
\newcommand{\AROUTE}	{\bar{\mathsf{R}}}
\newcommand{\CANMOVE}	{\mathsf{canmove}}	
\newcommand{\OCCUPIED}	{\mathsf{occupied}}

\newcommand{\CX}[1]		{& {\LG{\left [ #1\right ]}}}
\newcommand{\CXT}[1]	{& \LG{\left [ \text{#1}\right ]}}

\newcommand{\RULE}[2] 	{
\frac{
		\begin{array}{@{}l@{~~~}l}
			#1
		\end{array}
	}{
		\begin{array}{@{}l}
			#2
		\end{array}
	}
}



\definecolor{orange}{rgb}{0.9,1.0,0.9}
\definecolor{dgreen}{rgb}{0.0,0.4,0.0}
\definecolor{dblue}{rgb}{0.0,0.0,0.6}
\definecolor{lgray}{rgb}{0.4,0.4,0.4}
\definecolor{gray}{rgb}{0.3,0.3,0.3}

\newtheorem*{remark}{PhD Objective.}



\usetheme[progressbar=frametitle]{metropolis}
\usepackage{appendixnumberbeamer}

\usepackage{booktabs}
\usepackage[scale=2]{ccicons}

\usepackage{pgfplots}
\usepgfplotslibrary{dateplot}

\newcommand\blfootnote[1]{%
	\begingroup
	\renewcommand\thefootnote{}\footnote{#1}%
	\addtocounter{footnote}{-1}%
	\endgroup
}

\usepackage{xspace}
\newcommand{\themename}{\textbf{\textsc{metropolis}}\xspace}


\definecolor{darkblue}{rgb}{0,0,0.3}
\definecolor{lightblue}{rgb}{0.3,0.3,1}
\definecolor{darkred}{rgb}{0.7,0,0}
\definecolor{lightred}{rgb}{1,0.3,0.3}
\definecolor{orange}{rgb}{0.6,0.1,0.1}
\newcommand{\VAR}[1] {\textcolor{darkblue}{\mathit{#1}}}

\newcommand{\PAR}[1] {\textcolor{lightblue}{\mathit{#1}}}
\newcommand{\CST}[1] {\textcolor{lightred}{\mathsf{#1}}}

\newcommand{\fMODE} 			{\mathrm{MODE}}
\newcommand{\fMORDER} 			{\mathrm{ORDER}}
\newcommand{\fERROR} 			{\mathrm{ERROR}}
\newcommand{\fNOERROR} 			{\mathrm{NoError}}
\newcommand{\fINITMODE} 		{\mathrm{InitMode}}
\newcommand{\fMode}[1] 			{\mathrm{#1}}
\newcommand{\fConst}[1] 		{\mathrm{#1}}

\newcommand{\kG}[1]					{{\color{blue}{\texttt{#1}}}}
\newcommand{\kB}[1]					{\mathsf{#1}}
\newcommand{\kW}[1]					{{\footnotesize{\kG{#1}}}}

\newcommand{\kS}[1]					{{\footnotesize{$#1$}}}


\newcommand{\kwSYSTEM}			{machine}
\newcommand{\kwINTERFACE}		{INTERFACE}
\newcommand{\kwREFINEMENT}		{refinement}
\newcommand{\kwVARIABLES}		{variables}
\newcommand{\kwINVARIANT}		{invariant}
\newcommand{\kwINITIALISATION}	{initialisation}
\newcommand{\kwEVENTS}			{events}
\newcommand{\kwOPERATIONS}		{OPERATIONS}
\newcommand{\kwPROCESS}		        {PROCESS}
\newcommand{\kwMEND}				{end}
\newcommand{\kwPlink}		        {process link}
\newcommand{\kwWITH}		        {with}
\newcommand{\kwPARAMETERS}		{parameters}
\newcommand{\kwPROPERTIES}		{properties}
\newcommand{\kwANY}				{any}
\newcommand{\kwWHERE}			{where}
\newcommand{\kwTHEN}			{then}
\newcommand{\kwWHEN}			{when}
\newcommand{\kwSEES}			{sees}
\newcommand{\kwSKIP}			{skip}
\newcommand{\kwEND}				{end}
\newcommand{\kwBEGIN}			{begin}
\newcommand{\kwPRE}				{pre}
\newcommand{\kwPOST}			{post}
\newcommand{\kwREFINES}			{refines}
\newcommand{\LABEL}[1]			{\mathrm{#1}:~}

\newcommand{\Ztext}[1]			{\quad {\color{gray}{/\!/ ~ \text{#1}}}}

\newcommand{\EventBSnippet}[1] {
	\begin{center}
	\centerline{{\footnotesize{
	$
	\begin{array}{@{}l} #1 \end{array}
	$
	}}}
	\end{center}
}

\newcommand{\Snippet}[1] {
	\begin{center}
	\centerline{{\footnotesize{
	$
	\begin{array}{@{}l} #1 \end{array}
	$
	}}}
	\end{center}
}

\newcommand{\Eq}[2] {
	\begin{equation}
	\begin{array}{@{}l} #2 \end{array}
	\label{#1}
	\end{equation}	
}


\newcommand{\EventBSnippetBox}[1] {
	\centerline{{\fbox{\footnotesize{
	$
	\begin{array}{@{}l} #1 \end{array}
	$
	}}}}
}

\newcommand{\EventBSystem}[2] {
	\begin{center}
	{\normalsize{
	$
	\begin{array}{@{}l}
		\kG{\kwSYSTEM} ~ \kB{#1} \\
		\quad \begin{array}{@{}l} #2 \end{array} \\
		\kG{\kwMEND}
	\end{array}
	$
	}}
	\end{center}
}

\newcommand{\EventBSystemX}[2] {
	\begin{center}
	{\footnotesize{
	$
	\begin{array}{@{}l}
		\kG{\kwSYSTEM} ~ \kB{#1} \\
		\quad \begin{array}{@{}l} #2 \end{array} \\
	\end{array}
	$
	}}
	\end{center}
}


\newcommand{\EventBContext}[2] {
	{\footnotesize{
	$
	\begin{array}{@{}l}
		\kG{context} ~ \kB{#1} \\
		\quad \begin{array}{@{}l} #2 \end{array} \\
		\kG{\kwMEND}
	\end{array}
	$
	}}
}

\newcommand{\EventBContextd}[2] {
	{\footnotesize{
	$
	\begin{array}{@{}l}
		\kG{context} ~ \kB{#1} \\
		\quad \begin{array}{@{}l} #2 \end{array} \\
		%\kG{$\cdots$}
	\end{array}
	$
	}}
}


\newcommand{\EventBInterface}[2] {
	\begin{center}
	{\footnotesize{
	$
	\begin{array}{@{}l}
		\kG{\kwINTERFACE} ~ \kB{#1} \\
		\quad \begin{array}{@{}l} #2 \end{array} \\
		\kG{\kwMEND}
	\end{array}
	$
	}}
	\end{center}
}

\newcommand{\EventBRefinement}[2] {
	\begin{center}
	{\footnotesize{
	$
	\begin{array}{@{}l}
		\kG{\kwREFINEMENT} ~ \kB{#1} \\
		\quad \begin{array}{@{}l} #2 \end{array} \\
		\kG{\kwMEND}
	\end{array}
	$
	}}
	\end{center}
}




\newcommand{\EventBTemplate}[2] {
	$
	\begin{array}{@{}l}
		\kG{\kwTEMPLATE} ~ \kB{#1} \\
		\begin{array}{@{}l} #2 \end{array}
	\end{array}
	$
}


\newcommand{\EventBTemplateB}[2] {
	$\left (
	\begin{array}{@{}l}
		\kG{\kwTEMPLATE} ~ \kB{#1} \\
		\begin{array}{@{}l} #2 \end{array}
	\end{array}
	\right )$
}

\newcommand{\EventBGroup}[2] {


		\kG{group} ~ \kB{#1} ~ \kG{begin}\\
		\begin{array}{@{}l} \quad \begin{array}{rll} #2 \end{array} \end{array} \\
		\kG{end} \\
}

\newcommand{\EventBCont}[1] {
		\begin{center}
	{\footnotesize{
	$
	\begin{array}{@{~~~}l} #1 \end{array}
	$}}
	\end{center}
}

\newcommand{\EventBContX}[1] {
		\begin{center}
	{\footnotesize{
	$
	\begin{array}{@{}l}
		\begin{array}{@{}l} \quad \begin{array}{rll} #1 \end{array} \\ \kG{\kwMEND} \end{array}
	\end{array}
	$}}
		\end{center}
}


		


\newcommand{\EventBE}[1] {
	\begin{array}{@{}l} \quad \begin{array}{lll} #1 \end{array} \end{array}
}


\newcommand{\EventBRefines}[1] {
	\kG{\kwREFINES} ~ \kB{#1} \\
}

\newcommand{\EventBPlink} {
	\kG{\kwPlink} ~ \kB \\
}

\newcommand{\EventBImplements}[1] {
	\kG{implements} ~ \kB{#1} \\
}


\newcommand{\EventBExtends}[1] {
	\kG{extends} ~ \kB{#1} \\
}

\newcommand{\EventBSees}[1] {
	\kG{\kwSEES} ~ \kB{#1} \\
}


\newcommand{\EventBVars}[1] {
	\kG{\kwVARIABLES} \\
	\quad \begin{array}{l} #1 \end{array} \\
}

\newcommand{\EventBConstantsC}[1] {
	\kG{constants} ~ #1 \\
}

\newcommand{\EventBAxioms}[1] {
	\kG{axioms} \\
	\quad \begin{array}{l} #1 \end{array} \\
}

\newcommand{\EventBFlow}[1] {
	\kG{\kwFLOW} \\
	\quad \begin{array}{l} #1 \end{array} \\
}

\newcommand{\EventBVarsC}[1] {
	\kG{\kwVARIABLES} ~ #1 \\
}

\newcommand{\EventBAuxVarsC}[1] {
	\kG{auxiliary} ~ #1 \\
}

%\newcommand{\EventBConstantsC}[1] {
%	\kG{constants} ~ #1 \\
%}

%\newcommand{\EventBAxioms}[1] {
%	\kG{axioms} ~ #1 \\
%}

\newcommand{\EventBParamsC}[1] {
	\kG{\kwPARAMETERS} ~ #1 \\
}

\newcommand{\EventBProperties}[1] {
	\kG{\kwPROPERTIES} \\
	\quad \begin{array}{l} #1 \end{array} \\
}

\newcommand{\EventBInv}[1] {
	\kG{\kwINVARIANT} \\
	\quad \begin{array}{l} #1 \end{array} \\
}

\newcommand{\EventBVariant}[1] {
	\kG{variant} ~ #1 \\
}

\newcommand{\EventBInvC}[1] {
	\kG{\kwINVARIANT} ~ #1 \\
}

\newcommand{\EventBFaultsC}[1] {
	\kG{faults} ~ #1 \\
}

\newcommand{\EventBUsesC}[1] {
	\kG{uses} ~ \kB{#1} \\
}

\newcommand{\EventBUses}[2] {
	\kG{uses} ~ \kB{#1} \\
	\quad \begin{array}{l} #2 \end{array} \\
}


\newcommand{\EventBFaults}[1] {
	\kG{faults} \\
	\quad \begin{array}{l} #1 \end{array} \\
}

\newcommand{\Fault}[2] {
	\kB{#1} : #2 \\
}



\newcommand{\EventBInit}[1] {
	\kG{\kwINITIALISATION} \\
	\quad \begin{array}{l} #1 \end{array} \\
}

\newcommand{\EventBInitC}[1] {
	\kG{\kwINITIALISATION} ~ #1 \\
}

\newcommand{\EventBEvents}[1] {
	\kG{\kwEVENTS} \\
	\quad \begin{array}{lll} #1 \end{array} \\
}

%\newcommand{\EventBProcess}[1] {
%	\kG{\kwPROCESS} \\
%	\quad \begin{array}{lll} #1 \end{array} \\
%}



\newcommand{\EventBOperations}[1] {
	\kG{\kwOPERATIONS} \\
	\quad \begin{array}{lll} #1 \end{array} \\
}

\newcommand{\EventBProcess}[1] {
	\kG{\kwPROCESS} \\
	\quad \begin{array}{lll} #1 \end{array} \\
}


\newcommand{\EventBEventsTwoCol}[1] {
	\kG{\kwEVENTS} \\
	\quad \begin{array}{lll} #1 \end{array} \\
}

\newcommand{\EventBContTwoCol}[1] {
	$\begin{array}{lll} #1 \end{array}$
}

\newcommand{\EventBEventsRow}[2] {
	\begin{array}{lll} #1 \end{array} & \quad & \begin{array}{lll} #2 \end{array}\\
}

\newcommand{\EventBEventsRowSpan}[1] {
	\multicolumn{3}{l}{ \begin{array}{lll} #1 \end{array} }\\
}

\newcommand{\EventBAny}[4] {
				\kB{#1} = \\
				\kG{\kwANY} ~ #2 ~ \kG{\kwWHERE} \\
				\quad \begin{array}{l} #3 \end{array} \\
				\kG{\kwTHEN} \\
				\quad \begin{array}{l} #4 \end{array} \\
				\kG{\kwEND} \\
}

\newcommand{\EventBAnyInline}[4] {
				\kB{#1}  = \\
  			  	\kG{\kwANY} ~ #2 ~ \kG{\kwWHERE} \\
			   	\quad \begin{array}{l} #3 \end{array} \\
			   	\kG{\kwTHEN} \\
			   	\quad \begin{array}{l} #4 \end{array} \\
			   	\kG{\kwEND} \\
}

\newcommand{\EventBAnyWith}[5] {
  \kB{#1} & = 	& 	\kG{\kwANY} ~ #2 ~ \kG{\kwWHERE} \\
			  && 	\quad \begin{array}{l} #3 \end{array} \\
			  && 	\kG{\kwWITH} \\
			  && 	\quad \begin{array}{l} #4 \end{array} \\
			  && 	\kG{\kwTHEN} \\
			  && 	\quad \begin{array}{l} #5 \end{array} \\
			  && 	\kG{\kwEND} \\
}


\newcommand{\EventBAnyOperation}[4] {
  \kB{#1} & = 	& 	\kG{\kwANY} ~ #2 ~ \kG{\kwPRE} \\
			  && 	\quad \begin{array}{l} #3 \end{array} \\
			  && 	\kG{\kwPOST} \\
			  && 	\quad \begin{array}{l} #4 \end{array} \\
			  && 	\kG{\kwEND} \\
}

\newcommand{\EventBAnyOperationC}[4] {
  \kB{#1} & = 	& 	\kG{\kwANY} ~ #2 ~ \kG{\kwPRE} ~#3~ \kG{\kwPOST} ~#4~ \kG{\kwEND} \\
}

\newcommand{\EventBAnyRefines}[5] {
  \kB{#1} ~ \kG{refines} ~  #2 & = 	& 	\kG{\kwANY} ~ #3 ~ \kG{\kwWHERE} \\
			  && 	\quad \begin{array}{l} #4 \end{array} \\
			  && 	\kG{\kwTHEN} \\
			  && 	\quad \begin{array}{l} #5 \end{array} \\
			  && 	\kG{\kwEND} \\
}

\newcommand{\EventBAnyRefinesWith}[6] {
  \kB{#1} ~ \kG{refines} ~  #2 & = 	& 	\kG{\kwANY} ~ #3 ~ \kG{\kwWHERE} \\
			  && 	\quad \begin{array}{l} #4 \end{array} \\
			  && 	\kG{\kwWITH} \\
			  && 	\quad \begin{array}{l} #5 \end{array} \\
			  && 	\kG{\kwTHEN} \\
			  && 	\quad \begin{array}{l} #6 \end{array} \\
			  && 	\kG{\kwEND} \\
}


\newcommand{\EventBAnyC}[4] {
   \kB{#1} & = 	& 	\kG{\kwANY} ~ #2 ~ \kG{\kwWHERE} ~ #3 ~ \kG{\kwTHEN} ~#4~ \kG{\kwEND} \\
}


\newcommand{\EventBWhen}[3] {
   \kB{#1} & = 	& \kG{\kwWHEN} \\ 
			  && 	\quad \begin{array}{l} #2 \end{array} \\
			  && 	\kG{\kwTHEN} \\
			  && 	\quad \begin{array}{l} #3 \end{array} \\
			  && 	\kG{\kwEND} \\
}

\newcommand{\EventBWhenWith}[4] {
   \kB{#1} & = 	& 	\kG{with} \\ 
   			  &&	\quad \begin{array}{l} #2 \end{array} \\
   			  && 	\kG{\kwWHEN} \\ 
			  && 	\quad \begin{array}{l} #3 \end{array} \\
			  && 	\kG{\kwTHEN} \\
			  && 	\quad \begin{array}{l} #4 \end{array} \\
			  && 	\kG{\kwEND} \\
}

\newcommand{\EventBWhenOperation}[3] {
   \kB{#1} & = 	& \kG{\kwPRE} \\ 
			  && 	\quad \begin{array}{l} #2 \end{array} \\
			  && 	\kG{\kwPOST} \\
			  && 	\quad \begin{array}{l} #3 \end{array} \\
			  && 	\kG{\kwEND} \\
}

\newcommand{\EventBWhenRefines}[4] {
  \kB{#1} ~ \kG{ref} ~ #2 & = 	& \kG{\kwWHEN} \\ 
			  && 	\quad \begin{array}{l} #3 \end{array} \\
			  && 	\kG{\kwTHEN} \\
			  && 	\quad \begin{array}{l} #4 \end{array} \\
			  && 	\kG{\kwEND} \\
}

\newcommand{\EventBWhenC}[3] {
   \kB{#1} & = 	& \kG{\kwWHEN} ~#2~ \kG{\kwTHEN} ~#3~ \kG{\kwEND} \\
}

\newcommand{\EventBWhenOperationC}[3] {
   \kB{#1} & = 	& \kG{pre} ~#2~ \kG{post} ~#3~ \kG{\kwEND} \\
}

\newcommand{\EventBBegin}[2] {
   \kB{#1} & = 	& 	\kG{\kwBEGIN} \\
			  && 	\quad \begin{array}{l} #2 \end{array} \\
			  && 	\kG{\kwEND} \\
}

\newcommand{\EventBBeginC}[2] {
   \kB{#1} & = 	& 	\kG{\kwBEGIN} ~ #2 ~ \kG{\kwEND} \\
}

\newcommand{\EventBSkip}[1] {
   \kB{#1} & = 	& 	\kG{\kwSKIP} \\
}

\newcommand{\EventBSF}[2] {
   \kB{#1} & = 	& 	#2 \\
}




\usepackage{listings}
\usepackage{color}

\definecolor{codegreen}{rgb}{0,0.6,0}
\definecolor{codegray}{rgb}{0.5,0.5,0.5}
\definecolor{codepurple}{rgb}{0.58,0,0.82}
\definecolor{backcolour}{rgb}{0.95,0.95,0.92}

\lstdefinestyle{mystyle}{
	backgroundcolor=\color{backcolour},   
	commentstyle=\color{codegreen},
	keywordstyle=\color{magenta},
	numberstyle=\tiny\color{codegray},
	stringstyle=\color{codepurple},
	basicstyle=\scriptsize,
	breakatwhitespace=false,         
	breaklines=true,                 
	captionpos=b,                    
	keepspaces=true,                 
	numbers=left,                    
	numbersep=5pt,                  
	showspaces=false,                
	showstringspaces=false,
	showtabs=false,                  
	tabsize=2
}

\lstset{style=mystyle}

\title{Theories, Techniques and Tools for Engineering
	Heterogeneous Railway Networks}


%\date{}
\author{Paulius Stankaitis and Alexei Iliasov}
\institute{Centre for Software Reliability, Newcastle University, UK \vspace{0.2cm}\\ \begin{center} \small
		RSSRail Conference '17\vspace{0.1cm}\\ November 16th, Pistoia
	\end{center}}
\titlegraphic{\hfill\includegraphics[height=0.8cm]{Newcastle_University_Logo.pdf}}
\date{}
\begin{document}

\section{Scenario 1}

\begin{frame}{Dining Philosophers}
	\begin{center}
		
		\begin{tikzpicture}[scale=.5,every node/.style={minimum size=1cm},on grid]
		
		\node (p0)[process]  at (4,2) {$p_0$};
		\node (p1)[process]  at (8,2) {$p_1$};
		\node (r0)[resource] at (2,-1) {$r_{0}$};		
		\node (r1)[resource] at (6,-1) {$r_{1}$};
		\node (r2)[resource] at (10,-1) {$r_{2}$};
		
		\draw[requested] (r0) -- (p0) node [pos=0.1,right,font=\footnotesize] {q$_0$};
		\draw[requested] (r1) -- (p0) node [pos=0.1,left,font=\footnotesize] {q$_0$};
		
		\draw[requested] (r1) -- (p1) node [pos=0.1,right,font=\footnotesize] {q$_1$};
		\draw[requested] (r2) -- (p1) node [pos=0.1,left,font=\footnotesize] {q$_1$};
		
		
		\end{tikzpicture}
	\end{center}	
	
	\begin{center}
		
		
		\begin{tabular}{ c c| c c |c c} 
			
			&	r$_0$ &  & r$_1$ & & r$_2$   \\  
			\hline 
			0 &	  q$_0$*    & 1  &        & 2 &	\\ 
			1 &		  	   & 2  & 	    & 3	&			 \\ 
			2 &	      	   & 3  &        & 4 &                  \\ 
			\vdots &	\vdots & \vdots  & \vdots  & \vdots\\ 
			n & & n & &n
			
		\end{tabular}
		
	\end{center}
\textbf{Comment}: Suppose r$_0$ replies to p$_0$ first with 0 and increments pp(r$_0$) to pp(r$_0$) = 1 (* means promised, not locked yet).
\end{frame}

\begin{frame}{Dining Philosophers}
	\begin{center}
		
		\begin{tikzpicture}[scale=.5,every node/.style={minimum size=1cm},on grid]
		
		\node (p0)[process]  at (4,2) {$p_0$};
		\node (p1)[process]  at (8,2) {$p_1$};
		\node (r0)[resource] at (2,-1) {$r_{0}$};		
		\node (r1)[resource] at (6,-1) {$r_{1}$};
		\node (r2)[resource] at (10,-1) {$r_{2}$};
		
		\draw[requested] (r0) -- (p0) node [pos=0.1,right,font=\footnotesize] {q$_0$};
		\draw[requested] (r1) -- (p0) node [pos=0.1,left,font=\footnotesize] {q$_0$};
		
		\draw[requested] (r1) -- (p1) node [pos=0.1,right,font=\footnotesize] {q$_1$};
		\draw[requested] (r2) -- (p1) node [pos=0.1,left,font=\footnotesize] {q$_1$};
		
		
		\end{tikzpicture}
	\end{center}	
	
	\begin{center}
		
		
		\begin{tabular}{ c c| c c |c c} 
			
			&	r$_0$ &  & r$_1$ & & r$_2$   \\  
			\hline 
			0 &	  q$_0$*    & 1  &  q$_1$*       & 2 &	\\ 
			1 &		  	   & 2  & 	    & 3	&			 \\ 
			2 &	      	   & 3  &        & 4 &                  \\ 
			\vdots &	\vdots & \vdots  & \vdots  & \vdots\\ 
			n & & n & &n
			
		\end{tabular}
		
	\end{center}
	\textbf{Comment}: Suppose r$_1$ replies to p$_1$ with 1 and increments pp(r$_1$) to pp(r$_1$) = 2.
\end{frame}

\begin{frame}{Dining Philosophers}
	\begin{center}
		
		\begin{tikzpicture}[scale=.5,every node/.style={minimum size=1cm},on grid]
		
		\node (p0)[process]  at (4,2) {$p_0$};
		\node (p1)[process]  at (8,2) {$p_1$};
		\node (r0)[resource] at (2,-1) {$r_{0}$};		
		\node (r1)[resource] at (6,-1) {$r_{1}$};
		\node (r2)[resource] at (10,-1) {$r_{2}$};
		
		\draw[requested] (r0) -- (p0) node [pos=0.1,right,font=\footnotesize] {q$_0$};
		\draw[requested] (r1) -- (p0) node [pos=0.1,left,font=\footnotesize] {q$_0$};
		
		\draw[requested] (r1) -- (p1) node [pos=0.1,right,font=\footnotesize] {q$_1$};
		\draw[requested] (r2) -- (p1) node [pos=0.1,left,font=\footnotesize] {q$_1$};
		
		
		\end{tikzpicture}
	\end{center}	
	
	\begin{center}
		
		
		\begin{tabular}{ c c| c c |c c} 
			
			&	r$_0$ &  & r$_1$ & & r$_2$   \\  
			\hline 
			0 &	  q$_0$*    & 1  &  q$_1$*       & 2 & q$_1$*	\\ 
			1 &		  	   & 2  & 	    & 3	& 			 \\ 
			2 &	      	   & 3  &        & 4 &                  \\ 
			\vdots &	\vdots & \vdots  & \vdots  & \vdots\\ 
			n & & n & &n
			
		\end{tabular}
		
	\end{center}
	\textbf{Comment}: Suppose r$_2$ replies to p$_1$ with 2 and increments pp(r$_2$) to pp(r$_1$) = 3.
\end{frame}

\begin{frame}{Dining Philosophers}
	\begin{center}
		
		\begin{tikzpicture}[scale=.5,every node/.style={minimum size=1cm},on grid]
		
		\node (p0)[process]  at (4,2) {$p_0$};
		\node (p1)[process]  at (8,2) {$p_1$};
		\node (r0)[resource] at (2,-1) {$r_{0}$};		
		\node (r1)[resource] at (6,-1) {$r_{1}$};
		\node (r2)[resource] at (10,-1) {$r_{2}$};
		
		\draw[requested] (r0) -- (p0) node [pos=0.1,right,font=\footnotesize] {q$_0$};
		\draw[requested] (r1) -- (p0) node [pos=0.1,left,font=\footnotesize] {q$_0$};
		
		\draw[requested] (r1) -- (p1) node [pos=0.1,right,font=\footnotesize] {q$_1$};
		\draw[requested] (r2) -- (p1) node [pos=0.1,left,font=\footnotesize] {q$_1$};
		
		
		\end{tikzpicture}
	\end{center}	
	
	\begin{center}
		
		
		\begin{tabular}{ c c| c c |c c} 
			
			&	r$_0$ &  & r$_1$ & & r$_2$   \\  
			\hline 
			0 &	  q$_0$*    & 1  &  q$_1$*       & 2 & q$_1$*	\\ 
			1 &		  	   & 2  & 	q$_0$*    & 3	& 			 \\ 
			2 &	      	   & 3  &        & 4 &                  \\ 
			\vdots &	\vdots & \vdots  & \vdots  & \vdots\\ 
			n & & n & &n
			
		\end{tabular}
		
	\end{center}
	\textbf{Comment}: Suppose r$_1$ replies to p$_0$ with 2 and increments pp(r$_2$) to pp(r$_1$) = 3.
\end{frame}

\begin{frame}{Dining Philosophers}
	\begin{center}
		
		\begin{tikzpicture}[scale=.5,every node/.style={minimum size=1cm},on grid]
		
		\node (p0)[process]  at (4,2) {$p_0$};
		\node (p1)[process]  at (8,2) {$p_1$};
		\node (r0)[resource] at (2,-1) {$r_{0}$};		
		\node (r1)[resource] at (6,-1) {$r_{1}$};
		\node (r2)[resource] at (10,-1) {$r_{2}$};
		
		\draw[requested] (r0) -- (p0) node [pos=0.1,right,font=\footnotesize] {q$_0$};
		\draw[requested] (r1) -- (p0) node [pos=0.1,left,font=\footnotesize] {q$_0$};
		
		\draw[requested] (r1) -- (p1) node [pos=0.1,right,font=\footnotesize] {q$_1$};
		\draw[requested] (r2) -- (p1) node [pos=0.1,left,font=\footnotesize] {q$_1$};
		
		
		\end{tikzpicture}
	\end{center}	
	
	\begin{center}
		
		
		\begin{tabular}{ c c| c c |c c} 
			
			&	r$_0$ &  & r$_1$ & & r$_2$   \\  
			\hline 
			0 &	  q$_0$*    & 1  &  q$_1$*       & 2 & q$_1$*	\\ 
			1 &		  	   & 2  & 	q$_0$*    & 3	& 			 \\ 
			2 &	      	   & 3  &        & 4 &                  \\ 
			\vdots &	\vdots & \vdots  & \vdots  & \vdots\\ 
			n & & n & &n
			
		\end{tabular}
		
	\end{center}
	\textbf{Comment}: Both p$_0$ and p$_1$ can confirm they received all replies they needed, but max(p$_0$) = max(p$_1$) = 2 and so they both would like to lock r$_1$(2).
\end{frame}

\section{Scenario 2}
 \begin{frame}{Dining Philosophers}
	\begin{center}
		
		\begin{tikzpicture}[scale=.5,every node/.style={minimum size=1cm},on grid]
		
		\node (p0)[process]  at (0,2) {$p_0$};
		\node (p1)[process]  at (4,2) {$p_1$};
		\node (p2)[process]  at (8,2) {$p_2$};
		\node (r0)[resource] at (0,-1) {$r_{0}$};		
		\node (r1)[resource] at (5,-1) {$r_{1}$};
		\node (r2)[resource] at (10,-1) {$r_{2}$};
		
		\draw[requested] (r0) -- (p0);
		\draw[requested] (r1) -- (p0);
		
		\draw[requested] (r2) -- (p1);
		\draw[requested] (r0) -- (p1);
		
		\draw[requested] (r2) -- (p2);
		\draw[requested] (r0) -- (p2);
		
		
		\end{tikzpicture}
	\end{center}	
	
	\begin{center}
		
		
		\begin{tabular}{ c c| c c |c c} 
			
			&	r$_0$ &  & r$_1$ & & r$_2$   \\  
			\hline 
			0 &	 q$_0$*     & 1  &        & 2 &	\\ 
			1 &		  	   & 2  & 	    & 3	&			 \\ 
			2 &	      	   & 3  &        & 4 &                  \\ 
			\vdots &	\vdots & \vdots  & \vdots  & \vdots\\ 
			n & & n & &n
			
		\end{tabular}
		
	\end{center}
	\textbf{Comment}: Suppose r$_0$ replies to p$_0$ with 0 and increments pp(r$_0$) to pp(r$_0$) = 1.
\end{frame}

\begin{frame}{Dining Philosophers}
	\begin{center}
		
		\begin{tikzpicture}[scale=.5,every node/.style={minimum size=1cm},on grid]
		
		\node (p0)[process]  at (0,2) {$p_0$};
		\node (p1)[process]  at (4,2) {$p_1$};
		\node (p2)[process]  at (8,2) {$p_2$};
		\node (r0)[resource] at (0,-1) {$r_{0}$};		
		\node (r1)[resource] at (5,-1) {$r_{1}$};
		\node (r2)[resource] at (10,-1) {$r_{2}$};
		
		\draw[requested] (r0) -- (p0);
		\draw[requested] (r1) -- (p0);
		
		\draw[requested] (r2) -- (p1);
		\draw[requested] (r0) -- (p1);
		
		\draw[requested] (r2) -- (p2);
		\draw[requested] (r0) -- (p2);
		
		
		\end{tikzpicture}
	\end{center}	
	
	\begin{center}
		
		
		\begin{tabular}{ c c| c c |c c} 
			
			&	r$_0$ &  & r$_1$ & & r$_2$   \\  
			\hline 
			0 &	 q$_0$*     & 1  &        & 2 &	\\ 
			1 &	q$_1$*	  	   & 2  & 	    & 3	&			 \\ 
			2 &	 p$_2$*     	   & 3  &        & 4 &                  \\ 
			3 &	     	   & 4  &        & 5 &                  \\ 
			\vdots &	\vdots & \vdots  & \vdots  & \vdots\\ 
			n & & n & &n
			
		\end{tabular}
		
	\end{center}
	\textbf{Comment}: Suppose r$_0$ replies to p$_1$ and p$_2$ (r$_1$ hasn't replied to p$_0$ yet) respectively with 1 and 2, now pp(r$_0$) = 3
\end{frame}

\begin{frame}{Dining Philosophers}
	\begin{center}
		
		\begin{tikzpicture}[scale=.5,every node/.style={minimum size=1cm},on grid]
		
		\node (p0)[process]  at (0,2) {$p_0$};
		\node (p1)[process]  at (4,2) {$p_1$};
		\node (p2)[process]  at (8,2) {$p_2$};
		\node (r0)[resource] at (0,-1) {$r_{0}$};		
		\node (r1)[resource] at (5,-1) {$r_{1}$};
		\node (r2)[resource] at (10,-1) {$r_{2}$};
		
		\draw[requested] (r0) -- (p0);
		\draw[requested] (r1) -- (p0);
		
		\draw[requested] (r2) -- (p1);
		\draw[requested] (r0) -- (p1);
		
		\draw[requested] (r2) -- (p2);
		\draw[requested] (r0) -- (p2);
		
		
		\end{tikzpicture}
	\end{center}	
	
	\begin{center}
		
		
		\begin{tabular}{ c c| c c |c c} 
			
			&	r$_0$ &  & r$_1$ & & r$_2$   \\  
			\hline 
			0 &	 q$_0$*     & 1  &   q$_0$*      & 2 &	\\ 
			1 &	q$_1$*	  	   & 2  & 	    & 3	&			 \\ 
			2 &	 q$_2$*     	   & 3  &        & 4 &                  \\ 
			3 &	     	   & 4  &        & 5 &                  \\ 
		
			
		\end{tabular}
		
	\end{center}
	\textbf{Comment}: Suppose r$_1$ now replies to p$_0$ with 1 and increments pp(r$_1$) and now  p$_0$ could lock max(0, 1) = 1, but  r$_0$(1) is promised for someone else. Sure, it's possible that these will be released as other values might be higher, but it's likely they won't. (I'm considering locking lane with keys(q$_n$), but that's not a particularly nice solution)
\end{frame}


\end{document}
